% !TEX spellcheck = en_US



% ~~~~~~~~~~~~~~~~~~~~~~~~~~~~~~~~~~~~~~~ V
% (NC) 2026 Haluk Bingol
% github.com/halukbingol/LaTeX-Templates
%
% Licensees may copy, distribute, display, and perform the work and make derivative works 
% and remixes based on it only for non-commercial purposes. 
% ~~~~~~~~~~~~~~~~~~~~~~~~~~~~~~~~~~~~~~~ A


% v20261029T1800

\documentclass[aspectratio=169,10pt,compress]{beamer}




% ~~~~~~~~~~~~~~~~~~~~~~~~~~~~~~~~~~~~~~~~~~~~~~~~~~~~~~~~~~~~~~~~~ V hblib
	\usepackage{../../hblib/hblibPresentation} % lib presentation: theme, i/N, \hSection, algorithm2e, ...
	\usepackage{../../hblib/hblibCodeoutput} % lib code + output: \lstcode, \lstterminal, ...
	\usepackage{../../hblib/hblibDatastructures} % lib data structures: \hString, \hStack, \hQueue, \hArray
% ~~~~~~~~~~~~~~~~~~~~~~~~~~~~~~~~~~~~~~~~~~~~~~~~~~~~~~~~~~~~~~~~~ A hblib




% +++++++++++++++++++++++++++++++++++++++ V specific
\DontPrintSemicolon % in algorithms
% +++++++++++++++++++++++++++++++++++++++ A specific



\title[Arrays 1D: Senior]{
	One-Dimensional Arrays\\
	Senior
}
\subtitle{
	\cisOOO\ -- Level 4: internals, security, production use
}
\author[Erkovan]{Sait Erkovan}

\institute[FCIS]{
	Faculty of Computer and Information Sciences\\
	Yeditepe University
}
\date{{\tiny \hbTimeStamp}}

\begin{document}




% ~~~~~~~~~~~~~~~~~~~~~~~~~~~~~~~~~~~~~~ V frm
\begin{frame}[t,fragile]
	\titlepage
\end{frame}
% ~~~~~~~~~~~~~~~~~~~~~~~~~~~~~~~~~~~~~~ A frm




% ~~~~~~~~~~~~~~~~~~~~~~~~~~~~~~~~~~~~~~ V frm
\begin{frame}[t,fragile] \frametitle{Outline}
	\begin{columns}[T]
		\begin{column}{0.48\textwidth}
			\tableofcontents[sections={1-4}, hideallsubsections]
		\end{column}
		\begin{column}{0.48\textwidth}
			\tableofcontents[sections={5-8}, hideallsubsections]
		\end{column}
	\end{columns}

	\vspace{1em}
	\small Assumes \hCode{array1D-3-junior}. Facts about the JVM below are for the HotSpot JVM
	of OpenJDK 21; other JVMs may differ.
\end{frame}
% ~~~~~~~~~~~~~~~~~~~~~~~~~~~~~~~~~~~~~~ A frm




% ~~~~~~~~~~~~~~~~~~~~~~~~~~~~~~~~~~~~~~ V
\hSection[JVM]{Arrays inside the JVM}
% ~~~~~~~~~~~~~~~~~~~~~~~~~~~~~~~~~~~~~~ A




% ~~~~~~~~~~~~~~~~~~~~~~~~~~~~~~~~~~~~~~ V frm
\begin{frame}[t,fragile] \frametitle{An array is an object with a made-up class}
	\lstcoderightstep[0.66][basicstyle=\ttfamily\scriptsize]{codeoutput/ArrayTypes.java}{Java}{codeoutput/ArrayTypes.txt}

	\footnotesize \uncover<2->{\hCode{[} means ``array of''; \hCode{I D Z C} are \hCode{int double boolean char};
	\hCode{L\dots;} is a class. This explains \hCode{[I@\dots} in the freshman deck.}
\end{frame}
% ~~~~~~~~~~~~~~~~~~~~~~~~~~~~~~~~~~~~~~ A frm




% ~~~~~~~~~~~~~~~~~~~~~~~~~~~~~~~~~~~~~~ V frm
\begin{frame}[t,fragile] \frametitle{What an int[3] looks like in memory}
	\begin{center}
		\begin{tikzpicture}[cell/.style={draw, minimum height=0.6cm, font=\scriptsize, inner sep=1pt}]
			\node[cell, minimum width=3.2cm, fill=colorUnused] (m) at (0,0) {mark word: 8 bytes};
			\node[cell, minimum width=1.6cm, fill=colorUnused, right=0pt of m] (k) {class: 4};
			\node[cell, minimum width=1.6cm, fill=colorHead!50, right=0pt of k] (l) {length: 4};
			\node[cell, minimum width=1.6cm, right=0pt of l] (e0) {[0]: 4};
			\node[cell, minimum width=1.6cm, right=0pt of e0] (e1) {[1]: 4};
			\node[cell, minimum width=1.6cm, right=0pt of e1] (e2) {[2]: 4};
			\node[cell, minimum width=0.8cm, dashed, right=0pt of e2] (p) {};
			\draw[decorate, decoration={brace, amplitude=4pt}] (m.north west) -- (l.north east)
				node[midway, above=5pt, font=\scriptsize] {header: 16 bytes};
			\draw[decorate, decoration={brace, amplitude=4pt}] (e0.north west) -- (e2.north east)
				node[midway, above=5pt, font=\scriptsize] {elements: $3 \times 4$ bytes};
			\node[font=\scriptsize, below=2pt of p] {padding};
		\end{tikzpicture}
	\end{center}

	\begin{itemize}\small
		\item Typical 64-bit HotSpot with compressed class pointers; objects are aligned to 8 bytes:
		      \hCode{int[3]} takes 32 bytes, \hCode{int[1000]} about 4\,016
		\item The \hEmph{length is stored in the object}: \hCode{a.length} is a read, not a count
		\item For large arrays the header does not matter; for millions of tiny arrays it does
		\item Measure instead of guessing: the OpenJDK tool \emph{JOL} (Java Object Layout) prints the layout
		      \hRemark{(not run here)}
	\end{itemize}
\end{frame}
% ~~~~~~~~~~~~~~~~~~~~~~~~~~~~~~~~~~~~~~ A frm




% ~~~~~~~~~~~~~~~~~~~~~~~~~~~~~~~~~~~~~~ V frm
\begin{frame}[t,fragile] \frametitle{Arrays have their own bytecodes}
	\begin{columns}[T]
		\begin{column}{0.50\textwidth}
			\lstcode[basicstyle=\ttfamily\scriptsize]{codeoutput/Bytecode/Sum.java}{Java}
		\end{column}
		\begin{column}{0.48\textwidth}
			\uncover<2->{\lstterminal[basicstyle=\ttfamily\scriptsize]{codeoutput/Bytecode.txt}}
		\end{column}
	\end{columns}



% +++++++++++++++++++++++++++++++++++++++ V tbl
%: -Tbl. bytecodes
\begin{table}
		\scriptsize
		\begin{tabular}{@{}L{0.24\textwidth}L{0.68\textwidth}@{}}
			\toprule
			\textbf{Instruction} & \textbf{Java} \\
			\midrule
			\hCode{newarray int}   & \hCode{new int[3]}; checks that the length is not negative \\
			\hCode{iastore}        & \hCode{a[0] = 7}; checks \hCode{null} and the bounds \\
			\hCode{arraylength}    & \hCode{a.length} \\
			\hCode{iaload}         & \hCode{a[i]}; checks \hCode{null} and the bounds \\
			\bottomrule
		\end{tabular}
\end{table}%
% +++++++++++++++++++++++++++++++++++++++ A tbl


	\scriptsize \hCode{javap} comes with the JDK. Windows Command Prompt: \hCode{findstr} instead of \hCode{grep}.
\end{frame}
% ~~~~~~~~~~~~~~~~~~~~~~~~~~~~~~~~~~~~~~ A frm




% ~~~~~~~~~~~~~~~~~~~~~~~~~~~~~~~~~~~~~~ V
\hSection[Safety]{Checks that keep arrays safe}
% ~~~~~~~~~~~~~~~~~~~~~~~~~~~~~~~~~~~~~~ A




% ~~~~~~~~~~~~~~~~~~~~~~~~~~~~~~~~~~~~~~ V frm
\begin{frame}[t,fragile] \frametitle{Every access is checked}
	\begin{columns}[T]
		\begin{column}{0.48\textwidth}
			\textbf{Java}
			\begin{itemize}\small
				\item Every \hCode{a[i]} checks $0 \le i <$ \hCode{a.length}: a bad index is an
				      exception, never a silent read of other memory
				\item The JIT compiler removes a check when it can \emph{prove} it always passes,
				      e.g.\ in \hCode{for (i = 0; i < a.length; i++)}; so the safety is nearly free
			\end{itemize}
		\end{column}
		\begin{column}{0.48\textwidth}
			\textbf{C and C++}
			\begin{itemize}\small
				\item \hCode{a[i]} is plain address arithmetic: $s + i \cdot b$, no check
				\item Reading or writing past the end is a \hDefined{buffer overflow}:
				      one of the oldest sources of security holes
				\item Example: \emph{Heartbleed} (OpenSSL, 2014, CVE-2014-0160) read past a buffer
				      and leaked server memory
			\end{itemize}
		\end{column}
	\end{columns}

	\vspace{0.5em}
	\begin{block}{The other checks}
		\small \hCode{null} (NullPointerException), negative length (NegativeArraySizeException),
		element type (ArrayStoreException, next slide).
	\end{block}
\end{frame}
% ~~~~~~~~~~~~~~~~~~~~~~~~~~~~~~~~~~~~~~ A frm




% ~~~~~~~~~~~~~~~~~~~~~~~~~~~~~~~~~~~~~~ V frm
\begin{frame}[t,fragile] \frametitle{Covariance and ArrayStoreException}
	\lstcodedownstep[basicstyle=\ttfamily\scriptsize][basicstyle=\ttfamily\scriptsize]{codeoutput/Covariance.java}{Java}{codeoutput/Covariance.txt}

	\begin{itemize}\small
		\item<2-> Arrays are \hDefined{covariant}: a \hCode{String[]} \emph{is an} \hCode{Object[]}, so line 8 compiles
		\item<2-> The array remembers its element type; every store into an object array is checked at run time
		\item<2-> Generic types (\cisOOT) chose the opposite rule: such mistakes stop at compile time
	\end{itemize}
\end{frame}
% ~~~~~~~~~~~~~~~~~~~~~~~~~~~~~~~~~~~~~~ A frm




% ~~~~~~~~~~~~~~~~~~~~~~~~~~~~~~~~~~~~~~ V
\hSection[Limits]{Limits and overflow}
% ~~~~~~~~~~~~~~~~~~~~~~~~~~~~~~~~~~~~~~ A




% ~~~~~~~~~~~~~~~~~~~~~~~~~~~~~~~~~~~~~~ V frm
\begin{frame}[t,fragile] \frametitle{How big can an array be?}
	\begin{columns}[T]
		\begin{column}{0.70\textwidth}
			\lstcode[basicstyle=\ttfamily\scriptsize]{codeoutput/HeapLimits/Huge.java}{Java}
		\end{column}
		\begin{column}{0.28\textwidth}
			\begin{itemize}\small
				\item The index is an \hCode{int}: at most $2^{31} - 1$ elements
				\item The heap is limited: \hCode{-Xmx64m} gives 64 MB
			\end{itemize}
		\end{column}
	\end{columns}

	\vspace{0.5em}
	\small \hQuestion{What happens with 10 million, 20 million, and $2^{31} - 1$ \hCode{int}s?}
\end{frame}
% ~~~~~~~~~~~~~~~~~~~~~~~~~~~~~~~~~~~~~~ A frm




% ~~~~~~~~~~~~~~~~~~~~~~~~~~~~~~~~~~~~~~ V frm
\begin{frame}[t,fragile] \frametitle{Two different out-of-memory errors}
	\lstterminal[basicstyle=\ttfamily\scriptsize]{codeoutput/HeapLimits.txt}

	\begin{itemize}\small
		\item 10 million \hCode{int}s are 40 MB: they fit in 64 MB; 20 million (80 MB) do not
		\item $2^{31} - 1$ elements are refused before the heap is even asked: a JVM limit
		\item Large data: process it in chunks, or keep it in a file
	\end{itemize}
\end{frame}
% ~~~~~~~~~~~~~~~~~~~~~~~~~~~~~~~~~~~~~~ A frm




% ~~~~~~~~~~~~~~~~~~~~~~~~~~~~~~~~~~~~~~ V frm
\begin{frame}[t,fragile] \frametitle{A length computed with int can overflow}
	\lstcodedownstep[basicstyle=\ttfamily\scriptsize, firstline=4, lastline=16, firstnumber=4][basicstyle=\ttfamily\scriptsize]{codeoutput/SizeOverflow.java}{Java}{codeoutput/SizeOverflow.txt}

	\footnotesize $6 \cdot 10^{8} \cdot 4 > 2^{31} - 1$ wraps to a negative number; other counts wrap to a \emph{positive} wrong length, with no error.
\end{frame}
% ~~~~~~~~~~~~~~~~~~~~~~~~~~~~~~~~~~~~~~ A frm




% ~~~~~~~~~~~~~~~~~~~~~~~~~~~~~~~~~~~~~~ V frm
\begin{frame}[t,fragile] \frametitle{The midpoint bug}
	\lstcoderightstep[0.66][basicstyle=\ttfamily\scriptsize][basicstyle=\ttfamily\scriptsize]{codeoutput/Midpoint.java}{Java}{codeoutput/Midpoint.txt}

	\begin{itemize}\small
		\item \hCode{(lo + hi) / 2} was in binary searches and merge sorts for decades, including the JDK's
		      own \hCode{Arrays.binarySearch}; it fails only for arrays of more than about $2^{30}$ elements
		\item Found and published by J.\ Bloch in 2006
		      \hUrl{research.google/blog: Nearly All Binary Searches and Mergesorts are Broken}
		\item Fix: \hCode{lo + (hi - lo) / 2} or \hCode{(lo + hi) >>> 1}
	\end{itemize}
\end{frame}
% ~~~~~~~~~~~~~~~~~~~~~~~~~~~~~~~~~~~~~~ A frm




% ~~~~~~~~~~~~~~~~~~~~~~~~~~~~~~~~~~~~~~ V
\hSection[Security]{Security}
% ~~~~~~~~~~~~~~~~~~~~~~~~~~~~~~~~~~~~~~ A




% ~~~~~~~~~~~~~~~~~~~~~~~~~~~~~~~~~~~~~~ V frm
\begin{frame}[t,fragile] \frametitle{Arrays are mutable: copy at the border}
	\lstcode[basicstyle=\ttfamily\scriptsize, lastline=19]{codeoutput/DefensiveCopy.java}{Java}

	\footnotesize A \hCode{private final int[]} field is not protected: whoever has the arrow can change the elements.
\end{frame}
% ~~~~~~~~~~~~~~~~~~~~~~~~~~~~~~~~~~~~~~ A frm




% ~~~~~~~~~~~~~~~~~~~~~~~~~~~~~~~~~~~~~~ V frm
\begin{frame}[t,fragile] \frametitle{Who can change what?}
	\lstcoderightstep[0.66][basicstyle=\ttfamily\scriptsize, firstline=21, firstnumber=21]{codeoutput/DefensiveCopy.java}{Java}{codeoutput/DefensiveCopy.txt}

	\footnotesize \uncover<2->{The copies protect \hCode{scores}; the leaky getter and the public array do not.}
\end{frame}
% ~~~~~~~~~~~~~~~~~~~~~~~~~~~~~~~~~~~~~~ A frm




% ~~~~~~~~~~~~~~~~~~~~~~~~~~~~~~~~~~~~~~ V frm
\begin{frame}[t,fragile] \frametitle{The rules behind it}


% +++++++++++++++++++++++++++++++++++++++ V tbl
%: -Tbl. cert rules
\begin{table}
		\footnotesize
		\begin{tabular}{@{}L{0.20\textwidth}L{0.72\textwidth}@{}}
			\toprule
			\textbf{SEI CERT} & \textbf{Rule (Oracle Coding Standard for Java)} \\
			\midrule
			OBJ05-J & Do not return references to private mutable class members \\
			OBJ06-J & Defensively copy mutable inputs and mutable internal components \\
			OBJ13-J & Ensure that references to mutable objects are not exposed (\hCode{public static final} arrays) \\
			\bottomrule
		\end{tabular}
\end{table}%
% +++++++++++++++++++++++++++++++++++++++ A tbl


	\begin{itemize}\small
		\item An array handed out is an array handed over: return \hCode{clone()}
		\item An array taken in may be changed later by the caller: store \hCode{clone()}
		\item Instead of \hCode{public static final String[] LEVELS}, a private array and a method that returns a copy
	\end{itemize}


\end{frame}
% ~~~~~~~~~~~~~~~~~~~~~~~~~~~~~~~~~~~~~~ A frm




% ~~~~~~~~~~~~~~~~~~~~~~~~~~~~~~~~~~~~~~ V frm
\begin{frame}[t,fragile] \frametitle{Secrets in a char[], wiped after use}
	\lstcoderightstep[0.66][basicstyle=\ttfamily\scriptsize]{codeoutput/ClearSecret.java}{Java}{codeoutput/ClearSecret.txt}

	\footnotesize A \hCode{String} cannot be wiped; that is why \hCode{java.io.Console.readPassword()} returns a \hCode{char[]}.
\end{frame}
% ~~~~~~~~~~~~~~~~~~~~~~~~~~~~~~~~~~~~~~ A frm




% ~~~~~~~~~~~~~~~~~~~~~~~~~~~~~~~~~~~~~~ V frm
\begin{frame}[t,fragile] \frametitle{Never trust an index from outside}
	\lstcodedownstep[basicstyle=\ttfamily\scriptsize, firstline=5, firstnumber=5][basicstyle=\ttfamily\fontsize{6.4}{7.6}\selectfont, linerange={1-2,7-8}]{codeoutput/CheckIndex.java}{Java}{codeoutput/CheckIndex.txt}

	\footnotesize \uncover<2->{(JDK-internal lines cut.) \hCode{Objects.checkIndex}, Java 9+, fails early, before anything uses the index.}
\end{frame}
% ~~~~~~~~~~~~~~~~~~~~~~~~~~~~~~~~~~~~~~ A frm




% ~~~~~~~~~~~~~~~~~~~~~~~~~~~~~~~~~~~~~~ V
\hSection[Production]{Production use}
% ~~~~~~~~~~~~~~~~~~~~~~~~~~~~~~~~~~~~~~ A




% ~~~~~~~~~~~~~~~~~~~~~~~~~~~~~~~~~~~~~~ V frm
\begin{frame}[t,fragile] \frametitle{Less known parts of java.util.Arrays}
	\lstcoderightstep[0.66][basicstyle=\ttfamily\scriptsize, firstline=4, firstnumber=4]{codeoutput/ArraysPro.java}{Java}{codeoutput/ArraysPro.txt}

	\footnotesize \hCode{mismatch} and \hCode{compare}: Java 9. \hCode{parallelSort}, \hCode{stream}: Java 8.
\end{frame}
% ~~~~~~~~~~~~~~~~~~~~~~~~~~~~~~~~~~~~~~ A frm




% ~~~~~~~~~~~~~~~~~~~~~~~~~~~~~~~~~~~~~~ V frm
\begin{frame}[t,fragile] \frametitle{Arrays in real systems}


% +++++++++++++++++++++++++++++++++++++++ V tbl
%: -Tbl. production
\begin{table}
		\footnotesize
		\begin{tabular}{@{}L{0.30\textwidth}L{0.62\textwidth}@{}}
			\toprule
			\textbf{Situation} & \textbf{Practice} \\
			\midrule
			Arrays as keys, in sets, in \hCode{equals}  & \hCode{equals} and \hCode{hashCode} of an array are by identity: use \hCode{Arrays.equals}, \hCode{Arrays.hashCode} \\
			Public API                                  & return copies, or read-only views; document who owns the array \\
			Numbers in bulk (signals, images, matrices) & primitive arrays: compact, cache friendly, no boxing \\
			Very large data                             & watch $2^{31} - 1$ and the heap; stream it in chunks \\
			Sorting objects                             & \hCode{Arrays.sort} is stable for objects; equal keys keep their order \\
			Many threads                                & do not share a plain array that is written; see the next section \\
			\bottomrule
		\end{tabular}
\end{table}%
% +++++++++++++++++++++++++++++++++++++++ A tbl


\end{frame}
% ~~~~~~~~~~~~~~~~~~~~~~~~~~~~~~~~~~~~~~ A frm




% ~~~~~~~~~~~~~~~~~~~~~~~~~~~~~~~~~~~~~~ V
\hSection[Threads]{Arrays and threads}
% ~~~~~~~~~~~~~~~~~~~~~~~~~~~~~~~~~~~~~~ A




% ~~~~~~~~~~~~~~~~~~~~~~~~~~~~~~~~~~~~~~ V frm
\begin{frame}[t,fragile] \frametitle{Lost updates}
	\lstcodedownstep[basicstyle=\ttfamily\scriptsize, firstline=5, lastline=17, firstnumber=5][basicstyle=\ttfamily\scriptsize]{codeoutput/LostUpdates.java}{Java}{codeoutput/LostUpdates.txt}

	\footnotesize \uncover<2->{The \hCode{int[]} total varies from run to run and machine to machine: two threads can read the same old value.}
\end{frame}
% ~~~~~~~~~~~~~~~~~~~~~~~~~~~~~~~~~~~~~~ A frm




% ~~~~~~~~~~~~~~~~~~~~~~~~~~~~~~~~~~~~~~ V frm
\begin{frame}[t,fragile] \frametitle{Sharing arrays between threads}


% +++++++++++++++++++++++++++++++++++++++ V tbl
%: -Tbl. threads
\begin{table}
		\footnotesize
		\begin{tabular}{@{}L{0.32\textwidth}L{0.60\textwidth}@{}}
			\toprule
			\textbf{Tool} & \textbf{What it guarantees} \\
			\midrule
			\hCode{AtomicIntegerArray}, \hCode{AtomicLongArray} & each element is updated atomically \\
			\hCode{synchronized} block around the update        & one thread at a time; slower under contention \\
			\hCode{volatile int[] a}                            & only the \emph{arrow} is volatile, not the elements \\
			one array per thread, combine at the end            & no sharing, no locks: often the fastest \\
			\hCode{Arrays.parallelSort}, parallel streams       & the library splits the work safely \\
			\bottomrule
		\end{tabular}
\end{table}%
% +++++++++++++++++++++++++++++++++++++++ A tbl


	\vspace{0.3em}
	\small A program with a data race may print the right answer on your laptop and the wrong one on a server.
	Reasoning, not testing, shows it is correct.
\end{frame}
% ~~~~~~~~~~~~~~~~~~~~~~~~~~~~~~~~~~~~~~ A frm




% ~~~~~~~~~~~~~~~~~~~~~~~~~~~~~~~~~~~~~~ V
\hSection[Projects]{Open-ended projects}
% ~~~~~~~~~~~~~~~~~~~~~~~~~~~~~~~~~~~~~~ A




% ~~~~~~~~~~~~~~~~~~~~~~~~~~~~~~~~~~~~~~ V frm
\begin{frame}[t,fragile] \frametitle{Projects}
	\begin{enumerate}\small
		\item \textbf{Ring buffer.} The last $n$ sensor readings in a fixed array with a moving start index;
		      no shifting. Test it with JUnit, including wrap-around.
		\item \textbf{Bit set.} Store 64 \hCode{boolean}s in one \hCode{long} of a \hCode{long[]}; implement
		      \hCode{set}, \hCode{get}, \hCode{count}. Compare its memory with \hCode{boolean[]}.
		\item \textbf{Sieve of Eratosthenes} up to $10^{8}$: \hCode{boolean[]} against your bit set; time and memory.
		\item \textbf{Image histogram.} A gray image as one \hCode{int[]} of pixels (row after row);
		      histogram equalization with a counting array and prefix sums.
		\item \textbf{Secure container.} A class that holds scores and never leaks its array;
		      write tests that \emph{try} to change it from outside.
		\item \textbf{Benchmark study.} Sequential against random access, \hCode{int[]} against \hCode{Integer[]},
		      with JMH; explain every number.
	\end{enumerate}
\end{frame}
% ~~~~~~~~~~~~~~~~~~~~~~~~~~~~~~~~~~~~~~ A frm




% ~~~~~~~~~~~~~~~~~~~~~~~~~~~~~~~~~~~~~~ V
\hSection[Summary]{Summary}
% ~~~~~~~~~~~~~~~~~~~~~~~~~~~~~~~~~~~~~~ A




% ~~~~~~~~~~~~~~~~~~~~~~~~~~~~~~~~~~~~~~ V frm
\begin{frame}[t,fragile] \frametitle{Cheat sheet}


% +++++++++++++++++++++++++++++++++++++++ V tbl
%: -Tbl. cheat sheet
\begin{table}
		\footnotesize
		\begin{tabular}{@{}L{0.26\textwidth}L{0.66\textwidth}@{}}
			\toprule
			\textbf{Topic} & \textbf{Remember} \\
			\midrule
			Representation   & an object: header with the length, then the elements side by side \\
			Class names      & \hCode{[I}, \hCode{[D}, \hCode{[Ljava.lang.String;}; superclass \hCode{Object} \\
			Checks           & bounds, \hCode{null}, negative length, store type: all at run time \\
			Limits           & at most $2^{31} - 1$ elements; the heap (\hCode{-Xmx}) \\
			Overflow         & \hCode{Math.multiplyExact} for sizes; \hCode{lo + (hi - lo) / 2} \\
			Security         & copy in, copy out; no \hCode{public static final} arrays; wipe secrets \\
			Untrusted index  & \hCode{Objects.checkIndex(i, a.length)} \\
			Threads          & \hCode{AtomicIntegerArray}, or one array per thread \\
			\bottomrule
		\end{tabular}
\end{table}%
% +++++++++++++++++++++++++++++++++++++++ A tbl


\end{frame}
% ~~~~~~~~~~~~~~~~~~~~~~~~~~~~~~~~~~~~~~ A frm




% ~~~~~~~~~~~~~~~~~~~~~~~~~~~~~~~~~~~~~~ V frm
\begin{frame}[t,fragile] \frametitle{References}
	\begin{itemize}\small
		\item The Java Language Specification, Chapter 10: Arrays \hUrl{docs.oracle.com/javase/specs}
		\item The Java Virtual Machine Specification, \hCode{newarray}, \hCode{iaload}, \hCode{arraylength}
		\item SEI CERT Oracle Coding Standard for Java: OBJ05-J, OBJ06-J, OBJ13-J \hUrl{wiki.sei.cmu.edu}
		\item J.\ Bloch, ``Extra, Extra -- Read All About It: Nearly All Binary Searches and Mergesorts are Broken'', 2006
		      \hUrl{research.google/blog}
		\item \hCode{java.util.Arrays}, JDK 21 API documentation; Knuth \cite{knuth1973artv1}
	\end{itemize}

	\vspace{0.5em}
	\setbeamertemplate{bibliography item}[text]
	\bibliographystyle{IEEEtran}
	\tiny\bibliography{../../hblib/hblibDataStructures}
\end{frame}
% ~~~~~~~~~~~~~~~~~~~~~~~~~~~~~~~~~~~~~~ A frm




% ~~~~~~~~~~~~~~~~~~~~~~~~~~~~~~~~~~~~~~ V frm
\begin{frame}[t,fragile]
	\begin{center}
		\vspace{4cm}
		\hRemark{\Huge Questions?}
	\end{center}
\end{frame}
% ~~~~~~~~~~~~~~~~~~~~~~~~~~~~~~~~~~~~~~ A frm



\end{document}
